% ============================================================
% Chapitre 8 : Regression Lineaire Simple
% ============================================================
\chapter{R\'egression Lin\'eaire Simple}
\label{ch:regression-simple}

\begin{center}
\textit{<<~La r\'egression est l'outil le plus utile et le plus mal utilis\'e de la statistique.~>>}
\end{center}

% ------------------------------------------------------------
\section{Exemple motivant}
% ------------------------------------------------------------

On dispose de donn\'ees sur la surface habitable (en m$^2$) et le prix de vente (en k\texteuro) de 15 appartements. On veut mod\'eliser la relation $\text{Prix}=f(\text{Surface})$ et pr\'edire le prix pour une surface donn\'ee.

% ------------------------------------------------------------
\section{Le mod\`ele}
% ------------------------------------------------------------

\begin{definition}[Mod\`ele de r\'egression lin\'eaire simple]
\index{regression lineaire@r\'egression lin\'eaire!simple}
\[
Y_i = \beta_0 + \beta_1 x_i + \varepsilon_i, \quad i=1,\ldots,n,
\]
avec les hypoth\`eses :
\begin{enumerate}
\item $\E[\varepsilon_i]=0$ (erreurs centr\'ees),
\item $\Var(\varepsilon_i)=\sigma^2$ (homoscédasticité),
\item $\Cov(\varepsilon_i,\varepsilon_j)=0$ pour $i\neq j$ (erreurs non corr\'el\'ees),
\item (optionnel) $\varepsilon_i\sim\mathcal{N}(0,\sigma^2)$ (normalit\'e).
\end{enumerate}
$\beta_0$ : ordonn\'ee \`a l'origine. $\beta_1$ : pente (effet marginal de $x$ sur $Y$).
\end{definition}

% ------------------------------------------------------------
\section{Estimation par moindres carr\'es ordinaires (MCO)}
% ------------------------------------------------------------

\begin{theoreme}[Estimateurs MCO]
\index{moindres carres@moindres carr\'es}
Les estimateurs qui minimisent $\sum_{i=1}^n (Y_i-\beta_0-\beta_1 x_i)^2$ sont :
\begin{align}
\hat{\beta}_1 &= \frac{\sum_{i=1}^n(x_i-\bar{x})(Y_i-\bar{Y})}{\sum_{i=1}^n(x_i-\bar{x})^2} = \frac{S_{xy}}{S_{xx}}, \label{eq:beta1}\\
\hat{\beta}_0 &= \bar{Y}-\hat{\beta}_1\bar{x}. \label{eq:beta0}
\end{align}
\end{theoreme}

\begin{proof}
Posons $Q(\beta_0,\beta_1)=\sum(Y_i-\beta_0-\beta_1 x_i)^2$. Les conditions du premier ordre sont :
\[
\frac{\partial Q}{\partial\beta_0}=-2\sum(Y_i-\beta_0-\beta_1 x_i)=0, \quad
\frac{\partial Q}{\partial\beta_1}=-2\sum x_i(Y_i-\beta_0-\beta_1 x_i)=0.
\]
La premi\`ere donne $\beta_0=\bar{Y}-\beta_1\bar{x}$. Substitution dans la seconde :
\[
\sum x_i(Y_i-\bar{Y}+\beta_1\bar{x}-\beta_1 x_i)=0 \implies \beta_1\sum(x_i-\bar{x})^2=\sum(x_i-\bar{x})(Y_i-\bar{Y}). \qedhere
\]
\end{proof}

\begin{theoreme}[Propri\'et\'es des MCO]
\label{thm:mco-props}
\begin{enumerate}
\item $\E[\hat{\beta}_1]=\beta_1$ et $\E[\hat{\beta}_0]=\beta_0$ (sans biais).
\item $\Var(\hat{\beta}_1)=\frac{\sigma^2}{S_{xx}}$ et $\Var(\hat{\beta}_0)=\sigma^2\left(\frac{1}{n}+\frac{\bar{x}^2}{S_{xx}}\right)$.
\item $\hat{\sigma}^2=\frac{1}{n-2}\sum_{i=1}^n(Y_i-\hat{Y}_i)^2=\frac{\mathrm{SCR}}{n-2}$ est sans biais pour $\sigma^2$.
\end{enumerate}
\end{theoreme}

\begin{theoreme}[Gauss--Markov]
\index{Gauss--Markov}
Sous les hypoth\`eses (1)--(3) du mod\`ele, les MCO sont les \textbf{meilleurs estimateurs lin\'eaires sans biais} (BLUE) : parmi tous les estimateurs lin\'eaires et sans biais, ils ont la plus petite variance.
\end{theoreme}

% ------------------------------------------------------------
\section{D\'ecomposition de la variance}
% ------------------------------------------------------------

\begin{theoreme}[D\'ecomposition SCT = SCE + SCR]
\index{decomposition de la variance@d\'ecomposition de la variance}
\[
\underbrace{\sum_{i=1}^n(Y_i-\bar{Y})^2}_{\mathrm{SCT}} = \underbrace{\sum_{i=1}^n(\hat{Y}_i-\bar{Y})^2}_{\mathrm{SCE}} + \underbrace{\sum_{i=1}^n(Y_i-\hat{Y}_i)^2}_{\mathrm{SCR}},
\]
o\`u SCT = somme des carr\'es totale, SCE = somme des carr\'es expliqu\'ee, SCR = somme des carr\'es r\'esiduelle.
\end{theoreme}

\begin{definition}[Coefficient de d\'etermination]
\index{R2@$R^2$}
\[
R^2 = \frac{\mathrm{SCE}}{\mathrm{SCT}} = 1-\frac{\mathrm{SCR}}{\mathrm{SCT}} \in[0,1].
\]
$R^2$ mesure la proportion de la variance de $Y$ expliqu\'ee par le mod\`ele. En r\'egression simple, $R^2=r_{xy}^2$.
\end{definition}

% ------------------------------------------------------------
\section{Inf\'erence sous normalit\'e}
% ------------------------------------------------------------

Sous $\varepsilon_i\iid\mathcal{N}(0,\sigma^2)$ :

\begin{theoreme}[Tests et IC pour les coefficients]
\[
\frac{\hat{\beta}_1-\beta_1}{\hat{\sigma}/\sqrt{S_{xx}}} \sim t(n-2), \qquad
\frac{\hat{\beta}_0-\beta_0}{\hat{\sigma}\sqrt{1/n+\bar{x}^2/S_{xx}}} \sim t(n-2).
\]
\begin{itemize}
\item \textbf{Test} $H_0:\beta_1=0$ (pas de relation lin\'eaire) : $T=\hat{\beta}_1/\mathrm{SE}(\hat{\beta}_1)\sim t(n-2)$.
\item \textbf{IC 95\%} pour $\beta_1$ : $\hat{\beta}_1\pm t_{0{,}025,n-2}\cdot\mathrm{SE}(\hat{\beta}_1)$.
\end{itemize}
\end{theoreme}

\begin{theoreme}[Test global de Fisher]
\[
F = \frac{\mathrm{SCE}/1}{\mathrm{SCR}/(n-2)} = \frac{R^2/(1)}{(1-R^2)/(n-2)} \sim F(1,n-2) \text{ sous } H_0:\beta_1=0.
\]
En r\'egression simple, $F=T^2$ o\`u $T$ est la statistique du test de $\beta_1$.
\end{theoreme}

\subsection{Pr\'ediction}

\begin{definition}[Intervalle de confiance vs.\ de pr\'ediction]
Pour $x=x_0$ :
\begin{align}
\text{IC pour } \E[Y|x_0] &: \hat{Y}_0\pm t_{\alpha/2,n-2}\cdot\hat{\sigma}\sqrt{\frac{1}{n}+\frac{(x_0-\bar{x})^2}{S_{xx}}}, \\
\text{IP pour } Y_{\mathrm{new}} &: \hat{Y}_0\pm t_{\alpha/2,n-2}\cdot\hat{\sigma}\sqrt{1+\frac{1}{n}+\frac{(x_0-\bar{x})^2}{S_{xx}}}.
\end{align}
L'intervalle de pr\'ediction est toujours plus large (incertitude additionnelle de $\varepsilon$).
\end{definition}

% ------------------------------------------------------------
\section{Diagnostic des r\'esidus}
% ------------------------------------------------------------

\begin{definition}[R\'esidus]
\index{residus@r\'esidus}
Les r\'esidus sont $e_i=Y_i-\hat{Y}_i$. On examine :
\begin{enumerate}
\item \textbf{R\'esidus vs.\ valeurs ajust\'ees} : d\'etecte non-lin\'earit\'e et h\'et\'erosc\'edasticit\'e.
\item \textbf{QQ-plot des r\'esidus} : v\'erifie la normalit\'e.
\item \textbf{R\'esidus vs.\ ordre} : d\'etecte l'autocorr\'elation.
\end{enumerate}
\end{definition}

\begin{center}
\begin{tikzpicture}
% Residual patterns
\begin{scope}
\node at (2,3.5) {\small Bon (homoscédastique)};
\draw[->] (0,0) -- (4,0) node[right] {\tiny $\hat{Y}$};
\draw[->] (0,-1.5) -- (0,1.5) node[above] {\tiny $e$};
\foreach \x/\y in {0.3/0.2, 0.7/-0.4, 1.0/0.5, 1.5/-0.3, 2.0/0.1,
  2.5/-0.5, 2.8/0.3, 3.0/-0.1, 3.3/0.4, 3.7/-0.2} {
  \fill (\x,\y) circle (1.5pt);
}
\end{scope}
\begin{scope}[xshift=6cm]
\node at (2,3.5) {\small Mauvais (h\'et\'erosc\'edastique)};
\draw[->] (0,0) -- (4,0) node[right] {\tiny $\hat{Y}$};
\draw[->] (0,-1.5) -- (0,1.5) node[above] {\tiny $e$};
\foreach \x/\y in {0.3/0.1, 0.7/-0.1, 1.0/0.2, 1.5/-0.3, 2.0/0.4,
  2.5/-0.6, 2.8/0.7, 3.0/-0.8, 3.3/1.0, 3.7/-1.1} {
  \fill (\x,\y) circle (1.5pt);
}
\end{scope}
\end{tikzpicture}
\end{center}

% ------------------------------------------------------------
\section{Impl\'ementation Python}
% ------------------------------------------------------------

\begin{python}
\begin{minted}{python}
import numpy as np
import matplotlib.pyplot as plt
from scipy import stats
import statsmodels.api as sm

# --- Donnees : surface et prix ---
surface = np.array([25, 30, 35, 40, 45, 50, 55, 60, 65, 70,
                    75, 80, 90, 100, 120])
prix = np.array([95, 110, 125, 140, 148, 160, 175, 185, 195, 210,
                 220, 240, 260, 290, 350])

# --- Regression avec scipy ---
slope, intercept, r_value, p_value, std_err = stats.linregress(surface, prix)
print(f"beta_1 = {slope:.4f} (SE={std_err:.4f}, p={p_value:.6f})")
print(f"beta_0 = {intercept:.4f}")
print(f"R^2 = {r_value**2:.4f}")

# --- Regression avec statsmodels (plus complet) ---
X = sm.add_constant(surface)
model = sm.OLS(prix, X).fit()
print(model.summary())

# --- Graphique de regression ---
fig, axes = plt.subplots(2, 2, figsize=(12, 10))

# 1. Droite de regression
x_pred = np.linspace(20, 130, 100)
y_pred = intercept + slope * x_pred
axes[0, 0].scatter(surface, prix, color='blue', label='Donnees')
axes[0, 0].plot(x_pred, y_pred, 'r-', lw=2, label='Regression')
axes[0, 0].set_xlabel('Surface (m$^2$)')
axes[0, 0].set_ylabel('Prix (kEUR)')
axes[0, 0].set_title('Regression lineaire simple')
axes[0, 0].legend()

# 2. Residus vs valeurs ajustees
y_hat = intercept + slope * surface
residuals = prix - y_hat
axes[0, 1].scatter(y_hat, residuals)
axes[0, 1].axhline(0, color='red', linestyle='--')
axes[0, 1].set_xlabel('Valeurs ajustees')
axes[0, 1].set_ylabel('Residus')
axes[0, 1].set_title('Residus vs Ajustees')

# 3. QQ-plot
stats.probplot(residuals, plot=axes[1, 0])
axes[1, 0].set_title('QQ-plot des residus')

# 4. Histogramme des residus
axes[1, 1].hist(residuals, bins=6, edgecolor='black', alpha=0.7)
axes[1, 1].set_title('Distribution des residus')

plt.tight_layout()
plt.savefig("ch08_regression.pdf")
plt.show()

# --- Prediction et intervalle ---
x_new = 85
predictions = model.get_prediction(sm.add_constant([x_new]))
pred_summary = predictions.summary_frame(alpha=0.05)
print(f"\nPrediction pour {x_new} m^2:")
print(pred_summary)
\end{minted}
\end{python}

% ------------------------------------------------------------
\section{Exercices}
% ------------------------------------------------------------

\begin{exercice}[$\star$ -- Calcul MCO]
Calculer \`a la main $\hat{\beta}_0$, $\hat{\beta}_1$ et $R^2$ pour les donn\'ees :
\begin{center}
\begin{tabular}{c|ccccc}
$x$ & 1 & 2 & 3 & 4 & 5 \\
\hline
$y$ & 2.1 & 3.8 & 6.2 & 7.9 & 10.3 \\
\end{tabular}
\end{center}
\end{exercice}

\begin{exercice}[$\star\star$ -- Gauss--Markov]
Montrer que parmi les estimateurs de la forme $\hat{\beta}_1=\sum c_i Y_i$ qui sont lin\'eaires et sans biais pour $\beta_1$, le MCO a la plus petite variance.
\end{exercice}

\begin{exercice}[$\star\star$ -- Diagnostics]
R\'ealiser une r\'egression de la consommation de carburant sur le poids du v\'ehicule (donn\'ees \texttt{mtcars}). Analyser les r\'esidus. La relation est-elle vraiment lin\'eaire ?
\end{exercice}

\begin{exercice}[$\star\star\star$ -- Projet]
Collecter des donn\'ees (immobilier, m\'et\'eo, etc.), ajuster un mod\`ele de r\'egression simple, tester $\beta_1=0$, analyser les r\'esidus, construire des intervalles de pr\'ediction. Rapport complet avec code Python.
\end{exercice}

% ------------------------------------------------------------
\begin{formules}
\begin{align*}
\hat{\beta}_1 &= \frac{S_{xy}}{S_{xx}}, \quad \hat{\beta}_0=\bar{Y}-\hat{\beta}_1\bar{x} \\
R^2 &= 1-\frac{\mathrm{SCR}}{\mathrm{SCT}} = r_{xy}^2 \\
T &= \frac{\hat{\beta}_1}{\hat{\sigma}/\sqrt{S_{xx}}} \sim t(n-2) \\
F &= \frac{\mathrm{SCE}/1}{\mathrm{SCR}/(n-2)} \sim F(1,n-2)
\end{align*}
\end{formules}
